\section{Exact modular normal forms at rational cusps}\label{sec:cusps}

The modular formula below is classical.  Its role here is to provide
an exact analytic representative, including phases, whose inverse
can then be studied without losing a flat term.
The eta transformation and the Dedekind-sum multiplier are given in
\cite[Equations 23.18.5--23.18.7]{dlmf-eta}; an elementary proof of the
transformation is available in \cite{kong-teo}.

\subsection{One factor with all phases specified}

For integers \(H,K\), \(K>0\), \((H,K)=1\), write
\[
 s(H,K)=\sum_{\ell=1}^{K-1}
 \frac{\ell}{K}
 \left(\frac{H\ell}{K}
       -\left\lfloor\frac{H\ell}{K}\right\rfloor-\frac12\right),
 \qquad s(H,1)=0 .
\]
Fix a cusp \(h/k\), and for each positive integer \(d\) put
\[
 g_d=(d,k),\qquad H_d=\frac{dh}{g_d},\qquad K_d=\frac{k}{g_d}.
\]
Choose \(a_d\) with \(a_dH_d\equiv-1\pmod{K_d}\), and define
\[
 \rho_d=\ee^{2\pi\ii a_d/K_d}.
\]
For \(K_d=1\), set \(\rho_d=1\).
The value of \(\rho_d\) is independent of the chosen representative
of \(a_d\).

\begin{lemma}[One-factor cusp identity]\label{lem:one-cusp}
If \(\tau=h/k+\ii t/(2\pi)\) and \(\Re t>0\), then
\begin{align}
 \eta(d\tau)
 ={}&\exp\left(\pi\ii\left[
          \frac{dh}{12k}-s(H_d,K_d)\right]\right)
 \left(\frac{2\pi g_d}{kd\,t}\right)^{1/2}\notag\\
 &\quad{}\times
 \exp\left(-\frac{\pi^2g_d^2}{6dk^2t}\right)
 P\left(\rho_d
       \exp\left(-\frac{4\pi^2g_d^2}{dk^2t}\right)\right).
 \label{eq:one-cusp}
\end{align}
The square root is the branch positive for \(t>0\).
\end{lemma}
\begin{proof}
Choose \(b_d\in\Z\) so that
\[
 \gamma_d=
 \begin{pmatrix}a_d&b_d\\K_d&-H_d\end{pmatrix}
 \quad\hbox{has determinant }1.
\]
A direct calculation gives
\[
 \gamma_d(d\tau)=\frac{a_d}{K_d}
                  +\frac{2\pi\ii}{K_d^2d\,t}.
\]
The eta multiplier, in the convention of \cite{dlmf-eta}, is
\[
 \epsilon(\gamma_d)=
 \exp\left(\pi\ii\left[
      \frac{a_d-H_d}{12K_d}+s(H_d,K_d)\right]\right),
\]
and
\[
 \eta(\gamma_d z)
   =\epsilon(\gamma_d)\,[-\ii(K_dz-H_d)]^{1/2}\eta(z).
\]
At \(z=d\tau\), the square root equals
\((K_ddt/(2\pi))^{1/2}\).
Insert the product definition of eta at \(\gamma_d(d\tau)\)
and solve for \(\eta(d\tau)\).  The \(a_d\) phase cancels, leaving
\(\pi\ii H_d/(12K_d)-\pi\ii s(H_d,K_d)\).
Since \(H_d/K_d=dh/k\), this is \cref{eq:one-cusp}.
All transformations take place in the upper half-plane:
\(\Re t>0\) implies \(\Re(1/t)>0\), so the dual product converges.
\end{proof}

\subsection{A finite product quotient}

Let \(\mathcal D\) be a finite set of distinct positive integers
with nonzero integer exponents \(r_d\).
Set
\begin{align}
 B&=\sum_{d\in\mathcal D}r_d,&
 D&=\sum_{d\in\mathcal D}dr_d,\notag\\
 A_{h,k}&=\frac{\pi^2}{6k^2}
         \sum_{d\in\mathcal D}\frac{r_dg_d^2}{d},&
 N&=\lcm\{d:d\in\mathcal D\},\label{eq:cusp-constants}\\
 S&=\frac{4\pi^2}{k^2N},&
 M_d&=\frac{Ng_d^2}{d},\qquad Q=\ee^{-S/t}.\notag
\end{align}
Each \(M_d\) is a positive integer.  Define
\begin{equation}\label{eq:R-general}
 R(Q)=\prod_{d\in\mathcal D}P(\rho_dQ^{M_d})^{r_d},
\end{equation}
and the nonzero constant
\begin{equation}\label{eq:kappa}
 \kappa_{h,k}=
 \prod_{d\in\mathcal D}
       \left(\frac{2\pi g_d}{kd}\right)^{r_d/2}
 \exp\left(-\pi\ii
       \sum_{d\in\mathcal D}r_ds(H_d,K_d)\right).
\end{equation}

\begin{theorem}[Exact finite-product cusp normal form]
\label{thm:cusp-normal-form}
For \(q=\ee^{2\pi\ii h/k-t}\), \(\Re t>0\),
\begin{equation}\label{eq:raw-cusp}
 \prod_{d\in\mathcal D}P(q^d)^{r_d}
 =\kappa_{h,k}\,t^{-B/2}
   \exp\left(-\frac{A_{h,k}}t+\frac{Dt}{24}\right)R(Q).
\end{equation}
For the eta quotient,
\begin{equation}\label{eq:eta-cusp}
 \prod_{d\in\mathcal D}\eta(d\tau)^{r_d}
 =\kappa_{h,k}\ee^{\pi\ii hD/(12k)}
    t^{-B/2}\exp\left(-\frac{A_{h,k}}t\right)R(Q).
\end{equation}
These are exact holomorphic identities, not finite asymptotic
expansions.
\end{theorem}
\begin{proof}
Multiply \cref{eq:one-cusp} with the indicated integer powers.
The equality \(g_d^2/d=M_d/N\) puts all exponentials on the
common grid \(Q^{M_d}\).
For the raw product, use
\[
 P(q^d)=\ee^{-\pi\ii d\tau/12}\eta(d\tau).
\]
This removes the phase \(\ee^{\pi\ii dh/(12k)}\) and introduces
\(\ee^{dt/24}\).  The finite products of the remaining constants
give \cref{eq:kappa}.
Negative integer powers cause no difficulty because each factor
is nonzero in its disk.
\end{proof}

At \(q=1\), the formula reduces to
\[
 P(\ee^{-dt})=
 \sqrt{\frac{2\pi}{dt}}\,
 \ee^{-\pi^2/(6dt)+dt/24}
 P(\ee^{-4\pi^2/(dt)}).
\]
The extension to a rational cusp is not just replacement of
\(t\) by \(k^2t\): the factors \(g_d\) and phases \(\rho_d\)
also change.

\subsection{Arithmetic coefficients and resonant cancellation}

The logarithm of the dual product is convergent for \(|Q|<1\):
\begin{align}
 L(Q)=\log R(Q)&=\sum_{n\ge1}\lambda_nQ^n,\notag\\
 \lambda_n&=
 -\sum_{\substack{d\in\mathcal D\\M_d\mid n}}
 r_d\,\sigma_{-1}(n/M_d)\,\rho_d^{n/M_d}.
 \label{eq:lambda-arithmetic}
\end{align}
This follows by expanding
\(\log P(z)=-\sum_{n\ge1}\sigma_{-1}(n)z^n\).
If \(R(Q)=\sum_{n\ge0}c_nQ^n\), then
\begin{equation}\label{eq:euler-recurrence}
 c_0=1,\qquad
 nc_n=\sum_{j=1}^n j\lambda_jc_{n-j}.
\end{equation}
Thus all coefficients belong to the cyclotomic number field generated
by the finitely many \(\rho_d\).
No numerical cancellation test is needed.

The first nonzero coefficient of \(R-1\) equals the first nonzero
coefficient of \(L\).  It must be found \emph{after} the finite sum
in \cref{eq:lambda-arithmetic}, because distinct factors can cancel
at the same action.
For arbitrary phase-product representations the cancellation can
even be complete.  For example, separating even and odd factors gives
\[
 P(-Q)=\frac{P(Q^2)^3}{P(Q)P(Q^4)}.
\]
The constant case must therefore be included whenever a general
phase-product input is allowed.

For \(0<r<1\),
\begin{equation}\label{eq:log-product-bound}
 \sup_{|Q|\le r}|\log R(Q)|
 \le \sum_{d\in\mathcal D}|r_d|\,
            \frac{r^{M_d}}{(1-r^{M_d})^2}.
\end{equation}
Indeed
\(\sum_{\ell,n\ge1}|z|^{\ell n}/\ell
 \le |z|/(1-|z|)^2\).
Together with Cauchy's formula, this supplies explicit coefficient
and tail estimates.

\subsection{The four cusp regimes}

Taking logarithms of \cref{eq:raw-cusp} and setting \(v=1/t\)
gives
\begin{equation}\label{eq:log-core-cusp}
 \log F=
 -A_{h,k}v+\frac B2\Log v+\frac{D}{24v}
 +\log\kappa_{h,k}+L(\ee^{-Sv}).
\end{equation}
This is the exact-core form of \cref{eq:elementary-core}.
A local logarithm of the nonzero target is understood.

\begin{center}
\begin{tabularx}{\textwidth}{@{}>{\raggedright\arraybackslash}p{.24\textwidth}>{\raggedright\arraybackslash}p{.28\textwidth}>{\raggedright\arraybackslash}X@{}}
\toprule
Condition & Perturbative behavior & Appropriate inverse coordinate\\
\midrule
\(A_{h,k}\ne0\) &
Exponential--power core &
\(v=1/t\), with a linear--logarithmic or Lambert core\\
\(A_{h,k}=0,\ B\ne0\) &
Power of \(t\) &
A chosen algebraic power branch, then exponential corrections\\
\(A_{h,k}=B=0,\ D\ne0\), raw products &
Nonconstant analytic background &
The finite-endpoint theorem after taking logarithms\\
\(A_{h,k}=B=D=0\), raw products &
Constant plus flat terms &
Puiseux inverse in \(Q\), then a logarithm\\
\bottomrule
\end{tabularx}
\end{center}

For eta quotients the \(D/(24v)\) term is absent, so
\(A_{h,k}=B=0\) already gives the last regime.
This difference is why a normalization cannot be omitted when
identifying a flat critical value.

For completeness, the equation
\(av+b\Log v=w-d_0\), when \(ab\ne0\), has the local expression
\[
 v=\frac ba\,W_j\left(
       \frac ab\,\ee^{(w-d_0)/b}\right),
\]
with the \(W\)-branch chosen to satisfy the original logarithm
equation.  Exponentiating an equation can otherwise introduce a
spurious sheet.  A nonzero \(c/v\) can be retained in the exact
core or treated by a separate analytic core correction.

The common nome in \cref{eq:cusp-constants} is convenient and need
not be primitive.  Replacing \(Q\) by \(Q^g\) changes a displayed
vanishing order.  The constructions below have nonzero consecutive
exponents \(m,m+1\), so their ramification cannot be removed by
such a rescaling.  All ramification statements elsewhere are
relative to the stated coordinate \(Q\).

Finally, for a fixed family of dilations the physical condition
\(Q\to0\) is exactly
\[
 \frac{\Re(1/t)}{k^2}\longrightarrow+\infty.
\]
On the real radial path this is \(k^2t\to0\).
A fixed-cusp expansion therefore does not imply a uniform
small-nome approximation when the denominator \(k\) grows as fast
as \(t^{-1/2}\).
